% Generated draft: requires mathematical and source review.
\documentclass[11pt]{article}
\usepackage[margin=25mm]{geometry}
\usepackage{fontspec}
\setmainfont{Latin Modern Roman}
\usepackage{amsmath,amssymb,mathtools,unicode-math}
\setmathfont{Latin Modern Math}
\usepackage{enumitem,needspace,xurl,hyperref}
\setlength{\parindent}{0pt}
\setlength{\parskip}{5pt}
\setlength{\emergencystretch}{3em}
\newcommand{\N}{\mathbb{N}}
\newcommand{\Z}{\mathbb{Z}}
\newcommand{\Q}{\mathbb{Q}}
\newcommand{\R}{\mathbb{R}}
\newcommand{\C}{\mathbb{C}}
\newcommand{\F}{\mathbb{F}}
\newcommand{\A}{\mathbb{A}}
\newcommand{\PP}{\mathbb{P}}
\newcommand{\E}{\mathbb{E}}
\newcommand{\eps}{\varepsilon}
\newcommand{\dd}{\,\mathrm{d}}
\begin{document}
\section*{UnsolvedMath ID: 30005547}
\textbf{Source problem code:} catalogue entry
\subsection{Definitions and mathematical statement}
A Cantor set is conventionally a nonempty compact, perfect, totally disconnected set. Here the ambient space is \(\C\), and an analytic map means a holomorphic function on an open subset of \(\C\).

The supplied formulation asks: for every pair of Cantor sets \(X,Y\subset\C\), are there at most countably many analytic maps defined on a neighborhood of \(X\) that send \(X\) into \(Y\)? For a fixed open neighborhood \(U\supseteq X\), the functions described by this wording form the set
\[
\mathcal H_U(X,Y)=\{f:U\to\C : f\text{ is holomorphic and }f(X)\subseteq Y\}.
\]
However, the source does not specify how maps with different neighborhood domains are to be counted: as functions with their domains, as germs along \(X\), or as restrictions to \(X\). A germ along \(X\) identifies two functions when they agree on some open neighborhood of \(X\) contained in both domains.

There is also an essential unresolved qualification concerning constant maps: every \(y\in Y\) gives a permitted function \(f_y(z)=y\) under the supplied wording. The record does not state a nonconstancy condition or another exclusion addressing these maps. No such restriction is imposed in this draft.

The original excerpt instead asks whether the analogue of an unspecified \emph{Theorem 1} holds for Cantor sets \(X\subseteq\C\) and analytic maps. That theorem is not supplied, so its hypotheses and conclusion cannot be recovered from this record. The excerpt additionally says that the preceding problem might become easier under the assumption \(\dim\mu>0\), and that even this remains open. Neither \(\mu\), its dimension, nor the preceding problem is defined here. Thus a precise intended countability conjecture requires recovery of the missing context; the displayed family records only the supplied clean wording.
\subsection{Short English statement}
The supplied question concerns countability of holomorphic maps carrying one complex Cantor set into another. Its intended formulation remains uncertain because constant maps are permitted by the wording, the convention for identifying maps is unspecified, and the referenced Theorem 1 is absent. The additional condition \(\dim\mu>0\) also lacks its necessary context.
\subsection{Sources}
\begin{itemize}
\item Local UnsolvedMath registry, \texttt{data/unsolvedmath/problems.json}, permanent ID \(30005547\), title \emph{Analytic Maps of Complex Cantor Sets}, source problem code \texttt{catalogue entry}. The supplied record lists the proposed year as 2023 and names no proposer.
\item \emph{Incidence Problems in Harmonic Analysis, Geometric Measure Theory, and Ergodic Theory} (2023), Oberwolfach Reports, pp. 1397--1452, DOI \texttt{10.4171/owr/2023/25}. The original excerpt points to p. 1425. Supplied report link: \url{https://doi.org/10.4171/owr/2023/25}.
\item Supplied literature source, labeled OWR 2023 25: \url{https://publications.mfo.de/bitstream/handle/mfo/4055/OWR_2023_25.pdf?sequence=1}.
\item Literature assessment dated exactly \texttt{2026-08-22}: the supplied curation labels the problem open and states that the report explicitly leaves the complex Cantor set and analytic map analogue open. This is supplied dated curation, not an independently verified literature assessment.
\item The registry labels the statement \texttt{corrected\_verified} and reports verification against the official report. These are supplied source-verification claims; they have not been independently checked for this draft and do not resolve the formulation issues above.
\end{itemize}
\end{document}
